% SPDX-License-Identifier: MIT
% Overleaf: select XeLaTeX and set main.tex as the main document.
% The bundled CityU artwork has separate terms: see NOTICE.md.
\documentclass[aspectratio=169,11pt,t]{beamer}

\usetheme{CityU}
\usepackage{booktabs}
\usetikzlibrary{arrows.meta,positioning}

\title[CityU Beamer · January 2025 identity]{A CityU presentation,\\typeset with \LaTeX}
\subtitle{Research talks · Teaching · Conference presentations}
\author{Your Name}
\institute{Department / Research Group\\City University of Hong Kong}
\date{\today}

\begin{document}

\cityutitlepage

\begin{frame}{Today's route}
  \begin{columns}[T,onlytextwidth]
    \begin{column}{0.48\textwidth}
      \tableofcontents
    \end{column}
    \begin{column}{0.48\textwidth}
      \begin{block}{A familiar identity, an editable deck}
        The four backgrounds come from the supplied January 2025 PPT.
        Everything you write remains native \LaTeX.
      \end{block}
      \small
      Start with \texttt{minimal.tex}; use this document as a reference.
    \end{column}
  \end{columns}
\end{frame}

\section{Start with the message}

\begin{frame}{Give the audience a clear message}
  \framesubtitle{Native Beamer text, lists, and emphasis}
  \begin{itemize}
    \item Lead with the question your audience should remember.
    \item Explain why the question matters in one sentence.
    \item Connect the method to the evidence.
    \item Use \alert{one accent colour} to guide attention.
  \end{itemize}
  \medskip
  \begin{block}{The takeaway}
    One slide should support one central idea.
  \end{block}
\end{frame}

\begin{frame}{Build an argument in two columns}
  \begin{columns}[T,onlytextwidth]
    \begin{column}{0.48\textwidth}
      \begin{block}{The question}
        What changes when we introduce a new method?
      \end{block}
      \begin{itemize}
        \item Define a meaningful baseline.
        \item Keep the comparison fair.
        \item State the evaluation criteria.
      \end{itemize}
    \end{column}
    \begin{column}{0.48\textwidth}
      \begin{exampleblock}{The evidence}
        Present the result and explain its practical meaning.
      \end{exampleblock}
      \begin{alertblock}{The limitation}
        Make the conditions and remaining uncertainty explicit.
      \end{alertblock}
    \end{column}
  \end{columns}
\end{frame}

\begin{frame}{Reveal ideas at a deliberate pace}
  \begin{enumerate}[<+->]
    \item Establish the problem.
    \item Introduce the approach.
    \item Explain the contribution.
  \end{enumerate}
  \bigskip
  \begin{block}{One frame, three PDF pages}
    Beamer overlays reveal content in place. The footer keeps the same
    frame number on all three steps.\cite{beamer-manual}
  \end{block}
\end{frame}

\section{Show the evidence}

\begin{frame}{Typeset the mathematics directly}
  \framesubtitle{An illustrative regularised least-squares objective}
  \begin{equation}
    \widehat{\boldsymbol{\theta}}
      = \operatorname*{arg\,min}_{\boldsymbol{\theta}\in\mathbb{R}^{d}}
      \left\{
        \frac{1}{n}\sum_{i=1}^{n}
        \bigl(y_i-\mathbf{x}_i^{\mathsf T}\boldsymbol{\theta}\bigr)^2
        + \lambda\lVert\boldsymbol{\theta}\rVert_2^2
      \right\}.
  \end{equation}
  \begin{itemize}
    \item $\mathbf{x}_i$ and $y_i$: an input vector and its response.
    \item $\lambda\geq0$: the regularisation parameter.
    \item $n$ and $d$: the number of observations and features.
  \end{itemize}
  \smallskip
  {\small Equations remain vector text, with native numbering and symbols.}
\end{frame}

\begin{frame}{Let a table support a precise claim}
  \framesubtitle{Illustrative values only; not experimental findings}
  \begin{table}
    \centering
    \renewcommand{\arraystretch}{1.2}
    \begin{tabular}{lrrr}
      \toprule
      Variant & Score $\uparrow$ & Time (s) $\downarrow$ & Parameters \\
      \midrule
      Baseline & 0.72 & 12.4 & 128 \\
      Variant A & 0.78 & 10.1 & 128 \\
      Variant B & \textbf{0.81} & \textbf{9.6} & 128 \\
      \bottomrule
    \end{tabular}
    \caption{A compact comparison using \texttt{booktabs}.}
  \end{table}
  \begin{block}{Explain the comparison}
    Replace the sample values, describe the protocol, and report uncertainty
    where it is relevant.
  \end{block}
\end{frame}

\begin{frame}{Draw a reproducible workflow}
  \begin{figure}
    \centering
    \begin{tikzpicture}[
      stage/.style={draw=CityULine,fill=CityUPaper,rounded corners=2pt,
        minimum height=1.35cm,text width=3.3cm,align=center,font=\small},
      flow/.style={-{Latex[length=2mm]},thick,CityURed},
      node distance=9mm
    ]
      \node[stage] (input) {\textbf{Prepare}\\Define and inspect inputs};
      \node[stage,right=of input] (method) {\textbf{Analyse}\\Apply a stated method};
      \node[stage,right=of method] (output) {\textbf{Evaluate}\\Interpret the outputs};
      \draw[flow] (input) -- (method);
      \draw[flow] (method) -- (output);
    \end{tikzpicture}
    \caption{A native TikZ diagram with editable labels and vector geometry.}
  \end{figure}
  \medskip
  \begin{itemize}
    \item Use a diagram when a relationship is easier to see than to describe.
    \item Keep labels short and explain what moves between stages.
  \end{itemize}
\end{frame}

\begin{frame}{Cite sources without extra build tools}
  \small
  \begin{thebibliography}{9}
    \bibitem{beamer-manual}
      The Beamer Project.
      \newblock \emph{The Beamer class: User guide}.
      \newblock \href{https://ctan.org/pkg/beamer}{ctan.org/pkg/beamer}.

    \bibitem{cityu-identity}
      City University of Hong Kong, Communications and Institutional Research Office.
      \newblock Corporate identity resources.
      \newblock \href{https://www.cityu.edu.hk/ciro/}{cityu.edu.hk/ciro}.
  \end{thebibliography}
  \medskip
  \begin{block}{About this adaptation}
    The artwork was extracted from the supplied 202501 PPT.
    This is a community adaptation, not an officially endorsed template.
  \end{block}
\end{frame}

\section{Make it your own}

\begin{frame}[fragile]{Start from a minimal source file}
  \small
\begin{verbatim}
\documentclass[aspectratio=169,11pt,t]{beamer}
\usetheme{CityU}
\title[Short title]{Your presentation title}
\author{Your Name}
\begin{document}
\cityutitlepage
\begin{frame}{One idea} Text. \end{frame}
\cityuclosing{Thank you}
\end{document}
\end{verbatim}
  \begin{block}{Keep the files together}
    Compile beside \texttt{beamerthemeCityU.sty} and the \texttt{assets/}
    directory, using XeLaTeX.
  \end{block}
\end{frame}

\begin{frame}{Small controls, predictable results}
  \begin{description}[\texttt{sectionpages=false}]
    \item[\texttt{sectionpages=false}] Disable automatic section dividers.
    \item[\texttt{footer=false}] Hide the footer on content frames.
    \item[\texttt{assetspath=assets}] Set the relative artwork directory.
  \end{description}
  \medskip
  \begin{block}{Chinese and bilingual presentations}
    Open \texttt{example-zh.tex}. It uses \texttt{ctex} with Fandol fonts,
    which are distributed with TeX Live.
  \end{block}
  \smallskip
  {\small Cover, section, and closing helper pages do not count towards the footer total.}
\end{frame}

\begin{frame}{Before you share the deck}
  \begin{enumerate}
    \item Replace the example title, author, and illustrative content.
    \item Compile with XeLaTeX and inspect every PDF page.
    \item Check the permissions for any figures and university branding.
  \end{enumerate}
  \medskip
  \begin{alertblock}{Overleaf import is not Gallery approval}
    Public Gallery submission of a university presentation template requires
    official recognition and appropriate branding permissions.
    See the publication checklist in the documentation.
  \end{alertblock}
\end{frame}

\cityuclosing[Questions and discussion]{Thank you}

\end{document}
